%%%%%%%%%%%%%%%%%%%%%%%%%%%%%%%%%%%%%%%%%%%%%%%%%%%%%%%%
% 第6章 考察
%%%%%%%%%%%%%%%%%%%%%%%%%%%%%%%%%%%%%%%%%%%%%%%%%%%%%%%%
\chapter{考察}
\label{chap:discussion}

本章では，第5章で示した実験結果に基づき，
提案手法の有効性と限界について多角的な視点から考察を行う．
また，動物行動認識の実用化に向けた課題と展望についても議論する．

\section{マルチモーダル統合の有効性}

\subsection{外観特徴と運動特徴の相補性}
実験結果より，外観特徴（RGB）と運動特徴（オプティカルフロー）は
互いに相補的な役割を果たすことが確認された．
ソースドメイン内（Animal Kingdom）では両者が同程度の性能を示したのに対し，
ターゲットドメイン（ホッキョクグマ監視映像）では
運動特徴が外観特徴を大きく上回る性能を達成した．

この結果は，行動認識における各モダリティの役割が
ドメインシフトの程度によって変化することを示唆している．
学習データと評価データが同一分布に属する場合，
外観特徴は動物の姿勢や周囲の文脈情報を効果的に捉え，行動識別に寄与する．
一方，ドメインシフトが大きい場合，外観特徴は背景や照明条件の違いに
過度に敏感となり，行動の本質的な情報よりも環境依存的な情報を
捉えてしまう傾向がある．

運動特徴は，被写体の色や模様といった外観情報に依存せず，
純粋な動きのパターンを符号化できるため，
ドメイン間の環境差に対して頑健である．
この特性は，動物種や撮影環境が多様な実世界応用において極めて重要であり，
本研究の動機付けの妥当性を裏付けている．

\subsection{Gated Fusion の適応的統合能力}
提案した Gated Fusion 機構は，
入力サンプルに応じて両モダリティの寄与度を動的に調整する能力を持つ．
この機構により，動的な行動（走る，泳ぐなど）では運動特徴を重視し，
静的な行動（立つ，座るなど）では外観特徴を重視するといった，
柔軟な特徴統合が可能となる．

実験では，Gated Fusion を用いた手法が
単純な連結や固定重み付けによる統合と比較して
一貫して高い性能を示すことが確認された．
特に，ターゲットドメインにおいて
RGB Only（ARI: 0.093）と Flow Only（ARI: 0.557）の性能差が大きい中で，
Gated Fusion（ARI: 0.615）が両者を効果的に統合できた点は，
適応的統合機構の有効性を示す重要な証拠である．


\section{種に依存しない表現の獲得}

\subsection{敵対的学習によるドメイン不変性}
本研究では，動物種をドメインとみなし，
敵対的学習により種に依存しない行動表現の獲得を試みた．
実験結果より，敵対的学習の導入により
ターゲットドメインにおける性能が大幅に向上（$\Delta$ARI: +0.127）したことは，
この戦略の有効性を実証している．

従来のドメイン適応研究では，
ソースドメインとターゲットドメインが事前に定義されており，
ターゲットドメインのラベルなしデータを利用して適応を行うことが一般的であった \cite{ganin2016domain}．
一方，本研究では，ターゲットドメイン（ホッキョクグマ）のデータを
学習時に一切使用せず，ソースドメイン内の種多様性のみを活用して
種不変表現を学習した．
この設定は，事前にターゲットドメインが特定できない
実世界のシナリオにより適合しており，実用性が高い．

\subsection{種不変性と行動識別性のトレードオフ}
敵対的学習を導入する際の重要な課題は，
種情報の除去と行動識別能力の維持のバランスである．
過度に種情報を抑制すると，
行動識別に有用な情報まで失われる可能性がある．

本研究では，トリプレット損失による行動識別性の維持と
敵対的損失による種情報の抑制を同時に最適化することで，
このトレードオフに対処した．
Optuna によるハイパーパラメータ最適化の結果，
敵対的損失の重み $\lambda$ が適切に調整され，
両目的のバランスが取れた表現が獲得されたと考えられる．

ソースドメイン（Animal Kingdom）における性能が
敵対的学習の導入により大きく低下しなかった点（$\Delta$ARI: +0.025）は，
行動識別能力が維持されていることを示している．


\section{動物園実環境への適用可能性}

\subsection{常同行動検出への応用}
実験結果より，提案手法は「常同行動」を
他の行動クラスから明確に分離できることが確認された．
常同行動は動物福祉の重要な指標であり \cite{mason1991stereotypies}，
その自動検出は動物園における飼育管理の効率化に直結する．

特に，常同行動と通常の歩行（Walking）は動作パターンが類似するにもかかわらず，
提案手法はこれらを比較的良好に分離できた．
これは，敵対的学習により種固有のバイアスが除去され，
行動の反復パターンや軌跡の特徴がより純粋に抽出されたためと考えられる．

今後は，クラスタリングベースの教師なし異常検知手法と組み合わせることで，
過去の行動パターンからの逸脱を自動的に検出するシステムへの
拡張が期待される．

\subsection{24時間連続モニタリングへの展望}
本研究で用いたホッキョクグマ監視映像は，
動物園の固定カメラで撮影された実環境データである．
提案手法がこのような低解像度・低フレームレートの映像に対しても
有効に機能したことは，実運用への適用可能性を示唆している．

現在の動物園では，飼育員による定時観察が行われているが，
夜間や人手が不足する時間帯における継続的なモニタリングは困難である．
提案手法をエッジデバイスに実装し，
24時間連続でリアルタイムに行動をモニタリングするシステムを構築することで，
異常行動の早期発見や，行動パターンの長期的な変化の追跡が可能となる．

\subsection{他の動物種への汎化}
本研究では，ホッキョクグマに加えてアジアゾウの監視映像に対しても
提案手法の評価を行い，汎化性能の多角的な検証を試みた．
その結果，提案手法は両データセットにおいて最も高い性能を達成したものの，
アジアゾウにおいてはホッキョクグマと比較して全体的に性能が低下した．

この性能差の主要な要因として，以下の2点が考えられる：

\paragraph{動物種固有の動作特性}
ゾウはホッキョクグマと比較して体が大きく，全体的な動作が緩慢である傾向がある．
そのため，オプティカルフローによる運動特徴の抽出において，
異なる行動間の動作の差異が相対的に小さく，識別が困難になった可能性がある．
また，ゾウ特有の長い鼻（trunk）の動きが，
行動カテゴリとは独立したノイズとして作用している可能性もある．
これは，局所的な動作パターンが全体の行動特徴を覆い隠す形で影響を与えていると考えられる．

\paragraph{行動カテゴリの境界曖昧性}
ゾウの行動は，ホッキョクグマと比較して状態遷移が連続的で境界が曖昧なものが多い．
例えば，「立っている」状態から「ゆっくり歩いている」状態への遷移は
明確な境界を持たず連続的であり，人間のアノテーターにとっても区別が困難な場合がある．
このような行動定義の曖昧さは，真のラベルに対するクラスタリング性能の評価を
困難にしている可能性がある．

以上の考察より，提案手法の枠組みは他の動物種にも適用可能であることが
確認されたものの，対象動物の特性やデータ収集環境に応じた
追加的な工夫が有効である可能性が示唆された．
特に，複数個体環境への対応は，
実用的な動物園環境での展開において重要な課題である．

また，Animal Kingdom データセットは850種以上の動物を含んでおり \cite{ng2022animal}，
これらを学習に用いることで，より広範な動物種に対応できる
汎用的な行動認識システムの構築が期待される．
ただし，動物群によって行動様式が大きく異なるため，
行動クラスの定義や特徴抽出の方法については
さらなる検討が必要である．


\section{限界と今後の課題}
本研究の限界として，以下の点が挙げられる：

\paragraph{基盤モデルを用いた背景ノイズのロバスト化}
前述の通り，水面の波紋や植生の揺れなど，動的な背景要素がオプティカルフローに混入し，行動認識性能を低下させる場合がある．
この問題に対する従来のアプローチとして，背景差分法などが挙げられるが，動的な自然環境下では限界がある．
今後の課題として，Segment Anything Model (SAM) \cite{kirillov2023segment} などの強力なゼロショットセグメンテーション能力を持つ基盤モデルを導入し，
前処理段階で動物領域（前景）のみを明示的に切り出すことで，背景ノイズの影響を根本的に排除するパイプラインの構築が有効である．
また，VideoMAE \cite{tong2022videomae} のようなマスク映像モデリングを用いた事前学習を行うことで，
背景のテクスチャではなく，対象の構造的な動きに焦点を当てた特徴表現を獲得させることも期待される．

\paragraph{カメラワークの違いによるドメインシフト}
本研究の学習データ（ドキュメンタリー映像）はカメラワーク（パン・チルト・ズーム）を含む動的な映像が主である一方，
評価データ（監視カメラ映像）は固定カメラによる静的な映像である．
このカメラ動作の有無は，特にオプティカルフロー特徴における背景の動きの分布に大きな差異を生じさせる．
本実験ではオプティカルフローが高い汎化性能を示したが，
これは背景が静止している監視カメラ映像において，動物の動きのみが顕著に検出されたためと推測される．
逆のケース（固定カメラで学習し，手持ちカメラ映像で評価する場合）や，
風による植生の揺れなど背景動作が激しい環境においては，
カメラの動きと対象の動きを明示的に分離する技術が必要となる．

\paragraph{Animal Kingdom データセットにおける困難事例}
ソースドメインとして大規模な Animal Kingdom データセットを用いたが，
その中には行動認識が極めて困難な動画も含まれており，これらが学習のノイズとなった可能性がある．
実際の「難しい動画」の例として，以下のようなケースが挙げられる．
\begin{itemize}
    \item \textbf{激しいカメラブレ}: 手持ちカメラでの撮影等により画面全体が激しく揺れ，オプティカルフローが動物の動きを捉えきれない事例．
    \item \textbf{重度の遮蔽（オクルージョン）}: 草むらや木陰に動物が隠れており，体の一部しか視認できない事例．
    \item \textbf{被写体の微小性}: 広角で撮影されており，画面内における動物の領域が極めて小さく，テクスチャ情報が欠落している事例．
\end{itemize}
図\ref{fig:ak_difficult_samples}に，これらの困難な動画のフレーム例を示す．
（※論文で動画を紹介する際は，このように特徴的なフレームを数枚並べた図（フィルムストリップ形式）を用いるのが一般的である．）

\begin{figure}[t]
  \centering
  % \includegraphics[width=\linewidth]{image/ak_difficult_samples.png}
  \fbox{\begin{minipage}{0.9\linewidth}
    \centering
    \vspace{2cm}
    [Place holder for difficult samples: Camera shake, Occlusion, Small object]
    \vspace{2cm}
  \end{minipage}}
  \caption{Animal Kingdom データセットにおける認識困難な動画例（手ブレ，オクルージョン，微小な被写体）}
  \label{fig:ak_difficult_samples}
\end{figure}

これらの低品質なサンプルは，モデルが誤った特徴（背景の揺れなど）を行動と関連付けて学習する原因となり得る．
しかしながら，こうした「行動認識に適さない動画」を厳格にフィルタリングして除外すると，
学習に使用できるデータ数が大幅に減少し，深層学習モデルの訓練に必要な規模を維持できなくなるというジレンマがある．
実際，本研究の前処理段階（第4章参照）においても，クリップ長や単一個体条件などの最低限のフィルタリングを行っただけで，
元のデータセットから多数の動画が除外された．
したがって，本分野のさらなる発展のためには，
ノイズの多いサンプルに対してもロバストに学習できる手法の開発に加え，
動物行動分類に特化した，より高品質かつ精細な大規模動画データセットの構築が不可欠であると言える．
既存のデータセットはWeb上の動画を収集したものが多く，行動解析を主目的として撮影されていないため，
今後は研究目的に合致した撮影基準（固定カメラ，高解像度，遮蔽なし等）を満たす
標準データセットの整備が強く望まれる．

\paragraph{時間情報集約の高度化と周期性解析}
本研究の提案手法は，常同行動を含む多様な動作を高精度に分類できており，
クリップ単位での行動認識においては十分な成果を上げている．
より詳細な行動解析に向けては，集約手法の改善が考えられる．
現在用いている Average Pooling は，全体的な傾向を捉えるため，持続的な行動に対してはロバストであるが，
一瞬の決定的な動きが平滑化されてしまう側面がある．
これに対し，Max Pooling を用いれば，クリップ内で最も特徴的な瞬間を強調できる可能性があるが，
ノイズの影響を受けやすくなるリスクもあるため，一概に優れているとは言えない．
したがって，対象とする行動の特性に応じて手法を選択することや，
Attention Pooling や NetVLAD などの学習可能な集約層を導入し，
行動クラスの識別に真に重要なフレームを適応的に重み付けして選択することが有効と考えられる．

さらに，常同行動の本質である「反復性」や「持続性」を捉えるためには，
単純な集約だけでなく，長期間の周期性解析が不可欠である．
現在の数秒単位の特徴抽出に加え，数分から数十分のスケールで
特徴量の時系列変化に対して自己相関やフーリエ変換を適用することで，
ペーシングなどのリズムを定量化できる．
また，TimeSformer \cite{bertasius2021space} や ViViT \cite{arnab2021vivit} などの Video Transformer アーキテクチャを導入し，
Self-Attention 機構によって長距離の時間的依存関係を学習することで，
局所的な動作（歩行）と大域的なパターン（反復）を階層的に統合した理解が可能になると期待される．

\paragraph{実運用における行動境界の曖昧性}
本研究では，30秒の動画クリップに対してスライディングウィンドウとプーリングを適用し，特徴を集約して分類を行っている．
動物園での実運用を想定した場合，この固定長のウィンドウ内には「歩行」から「静止」，「採食」など，複数の異なる行動が混在する可能性が高い．
現状の特徴集約手法では，ウィンドウ内で支配的な動きに分類結果が引きずられたり，複数の行動特徴が混ざり合うことで識別精度が低下する懸念がある．
この課題に対しては，行動の開始点と終了点を明示的に推定する時間的アクションローカライゼーション（Temporal Action Localization）技術の導入や，
オンラインで行動の切り替わりを検出する変化点検知（Change Point Detection）を組み合わせることで，
より細粒度かつ正確な行動認識が可能になると考えられる．
